\section{Introduction}\label{sec:intro}
Within the standard cosmological model, dark energy drives the late-time accelerated expansion, but its physical origin remains uncertain \citep{PeeblesRatra2003DarkEnergy,Li2011DarkEnergy}. Testing whether it is a cosmological constant or evolves with time requires measurements of the expansion history and the growth of structure \citep{Weinberg2013Acceleration}. The large-scale distribution of galaxies provides a way to investigate these questions.

Galaxy clustering constrains distances and structure growth through two-point statistics, baryon acoustic oscillations, redshift-space distortions and the Alcock--Paczy\'nski (AP) effect \citep{Tegmark2004SDSS,BlakeGlazebrook2003BAO,SeoEisenstein2003BAO,Eisenstein2005BAO,Kaiser1987,Ballinger1996AP,Alam2017BOSS}. These applications motivate complementary probes of the geometry and connectivity of large-scale structure \citep{Blake2011Growth,Blake2011AP,Beutler2017Anisotropic}.

The genus characterizes the topology of the large-scale density field, complementing the power spectrum and two-point correlation function \citep{Hamilton1986Topology,Mecke1994,SchmalzingBuchert1997Morphology}. At a fixed smoothing scale, a uniform positive rescaling of the density contrast preserves the genus curve expressed in volume-fraction thresholds \citep{Matsubara2003SmoothedFields,Wang2012TopologyGravity}, motivating its use in studies of redshift evolution. Topological measurements in the 2dF Galaxy Redshift Survey, SDSS luminous red galaxy sample and WiggleZ survey provide observational applications of these statistics \citep{Park2005SDSS,James2009Topology2dF,Gott2009LRGGenus,Blake2014WiggleZTopology}. Simulations also show that genus evolution depends on tracer bias and gravitational growth \citep{Park2005Evolution}. More recently, the four Minkowski functionals of BOSS galaxies have been analysed with a simulation-based model for cosmological inference \citep{Liu2025}.

\citet{ParkKim2010Ruler} proposed using large-scale topology as a cosmic standard ruler, motivating comparisons of the genus across redshifts. This approach was developed through mock-catalogue tests of the genus amplitude \citep{Appleby2018} and measurements of the expansion history from the two-dimensional genus of observed galaxies \citep{Appleby2021GenusExpansion}. These studies provide the main motivation for examining whether H\,I topology retains a useful redshift-dependent geometric response after foreground filtering.

H\,I intensity mapping traces large-scale structure through integrated 21-cm emission \citep{Battye2013IntensityMapping}, and its clustering depends on how H\,I populates dark matter haloes \citep{Villaescusa2018}. Topology has been proposed as a probe of the H\,I distribution with the Square Kilometre Array (SKA) \citep{Wang2015HITopologySKA}. Minkowski functionals have also been used to study the H\,I--halo mass relation in two-dimensional intensity maps with foreground filtering \citep{Spina2021HIHaloMF}.

An assumed cosmology is required to convert redshifts and angles into comoving positions. An incorrect distance--redshift relation introduces redshift-dependent AP distortions, which can alter the measured genus at a fixed smoothing scale \citep{AlcockPaczynski1979}. Following the redshift-differencing approach of tomographic AP analyses \citep{Li2014GradientAP,Li2016TomographicAP,Park2019EvolutionAP}, we compare the full three-dimensional genus curves of H\,I mass-weighted fields across redshifts after subtracting the intrinsic evolution measured in the simulation cosmology.

Foreground avoidance in 21-cm surveys removes contaminated Fourier modes at low line-of-sight wavenumbers and within the foreground wedge \citep{SeoHirata2016Wedge}. The wedge arises from mode mixing in the chromatic interferometer response \citep{Datta2010BrightSourceWedge,Morales2012ForegroundShapes}. We use idealized line-of-sight and wedge cuts to examine how mode removal affects the redshift evolution of the H\,I-weighted genus and its sensitivity to the assumed cosmology. The cuts are not calibrated to a specific instrument. Frequency-dependent beam errors can introduce additional foreground leakage \citep{Ding2024BeamCalibration}; this instrumental effect is not included here.

We address three questions: how foreground removal changes the genus and its evolution in the simulation cosmology; whether redshift-differenced genus measurements retain a response to an incorrect distance--redshift relation; and how reliably the resulting parameter regions can be compared. Sections~\ref{sec:data}--\ref{sec:statistics} describe the samples, measurement and statistic. Section~\ref{sec:results} presents the field morphology, genus evolution, geometric response and parameter regions, together with numerical sensitivity tests. Section~\ref{sec:discussion} discusses their implications for H\,I surveys.
